% !TEX root = ../../main.tex

% --- القسم الأول: توصيف المنظومة ونظرية الترشيح ---
\section{توصيف المنظومة ونظرية الترشيح}

تستكمل هذه الدراسة ما بدأناه في النمذجة السابقة، حيث نركز هنا على تطبيق نوع آخر من تقنيات الترشيح لتحسين جودة القدرة الكهربائية في الأنظمة التي تغذي أحمالاً غير خطية.

\begin{itemize}
    \item \textbf{هدف الوظيفة:} تصميم \textbf{مرشح تمرير منخفض \el{(Low Pass Filter)}} واختيار بارامتراته بدقة لتقليل نسبة التشوه التوافقي الكلي \el{($THD$)} للجهد والتيار، وضمان بقاء النتائج ضمن المعايير العالمية المسموح بها (أقل من \el{3\%})، وحماية عناصر القدرة من النبضات والضجيج.
    
    \item \textbf{مبدأ العمل:} يعمل مرشح التمرير المنخفض كأداة كهربائية تسمح بمرور التردد الأساسي ($50 Hz$) بسهولة نتيجة الممانعة المنخفضة للملف ($X_L$) والممانعة العالية للمكثف ($X_C$) عند هذا التردد. في المقابل، يفرض المرشح ممانعة عالية أمام التوافقيات ذات الترددات الكبيرة (مثل $250 Hz$ و $350 Hz$)، مما يؤدي إلى إخمادها وتبديدها قبل دخولها إلى المقوم.
    
    \item \textbf{توصيف الدارة:} لضمان اتساق الدراسة والمقارنة الدقيقة، سنحافظ على نفس مكونات الشبكة الموصوفة في القسم الأول من الوظيفة السابقة، والتي تشمل منبع الجهد ثلاثي الطور ($400 V, 50 Hz$)، ممانعة القصر ($Z_{sc}$)، ممانعة المحولة ($Z_T$)، وممانعة الكبل، وصولاً إلى الحمل غير الخطي المتمثل بجسر التقويم الديودي.
\end{itemize}
